% Generated draft: requires mathematical and source review.
\documentclass[11pt]{article}
\usepackage[margin=25mm]{geometry}
\usepackage{fontspec}
\setmainfont{Latin Modern Roman}
\usepackage{amsmath,amssymb,mathtools,unicode-math}
\setmathfont{Latin Modern Math}
\usepackage{enumitem,needspace,xurl,hyperref}
\setlength{\parindent}{0pt}
\setlength{\parskip}{5pt}
\setlength{\emergencystretch}{3em}
\newcommand{\N}{\mathbb{N}}
\newcommand{\Z}{\mathbb{Z}}
\newcommand{\Q}{\mathbb{Q}}
\newcommand{\R}{\mathbb{R}}
\newcommand{\C}{\mathbb{C}}
\newcommand{\F}{\mathbb{F}}
\newcommand{\A}{\mathbb{A}}
\newcommand{\PP}{\mathbb{P}}
\newcommand{\E}{\mathbb{E}}
\newcommand{\eps}{\varepsilon}
\newcommand{\dd}{\,\mathrm{d}}
\begin{document}
\section*{UnsolvedMath ID: 8800017}
\textbf{Source problem code:} catalogue entry
\subsection{Definitions and mathematical statement}
In the finite-field setting suggested by the workshop, let $S/\mathbb F_q$ be a surface in a three-dimensional toric variety. Its zeta function is
\[
Z(S,t)=\exp\left(\sum_{r\geq1}\frac{\#S(\mathbb F_{q^r})}{r}t^r\right).
\]
Develop effective algorithms computing this rational function exactly for the source's nondegenerate surfaces in toric threefolds.

A standard hypersurface setting starts with a Laurent polynomial $f\in\mathbb F_q[x_1^{\pm1},x_2^{\pm1},x_3^{\pm1}]$ and its Newton polytope $\Delta$, the convex hull of its exponent vectors. For a nonempty face $\tau\subseteq\Delta$, let $f_\tau$ be the sum of terms on that face. Newton nondegeneracy means that
\[
f_\tau=x_1\frac{\partial f_\tau}{\partial x_1}
=x_2\frac{\partial f_\tau}{\partial x_2}
=x_3\frac{\partial f_\tau}{\partial x_3}=0
\]
has no solution in $(\overline{\mathbb F}_q^*)^3$ for every such face, including $\Delta$. This gives a conventional meaning of nondegenerate, not an independently verified definition from the slide.

The input/output theorem should specify whether $S$ is the open torus hypersurface or its compactification in a toric threefold, the fan and smoothness hypotheses, and the exact required notion of effectiveness. The two zeta functions differ by boundary contributions. The source supplies no complexity target in $q$, characteristic, degrees, Newton-polytope size, or input length, so no polynomial-time bound is invented.
\subsection{Short English statement}
Compute exact zeta functions effectively for nondegenerate surfaces in toric threefolds over finite fields. Fix the nondegeneracy and compactification conventions and state the algorithm's input representation and cost guarantees.
\subsection{Sources}
\begin{itemize}
\item Kiran Kedlaya, IPAM \emph{Workshop on Number Theory and Cryptography --- Open Problems} (2006), attachment 6220, PDF pages 40--41, logical slide 21, schedule source \url{https://www.ipam.ucla.edu/programs/workshops/workshop-i-number-theory-and-cryptography-open-problems/?tab=schedule}; accessed 2026-07-29. The supplied provenance records a PDF/legacy-PPT visual audit.
\item The research summary mentions work attributed to Castryck--Chtcherbakov and Gajovic on toric/projective surfaces and higher-dimensional toric $p$-adic algorithms, but supplies no article titles, identifiers, or complete theorem hypotheses.
\item Local registry \texttt{data/unsolvedmath/problems.json}, permanent UnsolvedMath ID $8800017$, source code catalogue entry, \emph{Nondegenerate surfaces in toric threefolds}. Dataset attribution: Ulam AI and the UnsolvedMath Contributors.
\end{itemize}
\end{document}
